\section{Evaluation protocol}
\label{sec:protocol}

\subsection{Thresholds, and what ``fair'' means here}

Monitors live on incomparable scales --- a Frobenius norm and a cosine distance
have no common unit --- so comparing their raw scores is meaningless. What can
be matched is their false-alarm behaviour. For each monitor independently, the
threshold is the $99$th percentile of that monitor's scores pooled over the
\emph{control-calibration} clips and the evaluation frames. Every monitor
therefore raises false alarms on $1\%$ of control frames by construction, and
differences in detection latency are differences in sensitivity, not in scale.

Control clips are split into disjoint calibration ($n=\lgwNCalib$) and
evaluation ($n=\lgwNEval$) sets. Thresholds are set on the former and false
alarms measured on the latter; no threshold ever sees anomalous data.

\paragraph{A second operating point.}
Matching the \emph{per-frame} rate at $1\%$ compounds to roughly $14\%$ over a
$15$-frame episode, which no episode-level monitor could use. We therefore
repeat the entire analysis with thresholds re-derived so the \emph{clip-level}
false-alarm rate on calibration controls is $5\%$. This was fixed before any
inference run, and both regimes are reported unconditionally, so neither can be
selected after the fact.

\subsection{Detection latency as a survival quantity}

An alarm fires at the first evaluation frame whose score exceeds the threshold.
With injection at $t^\star = 15$ and an evaluation window of frames
$15..29$, detection latency is $\dt = t_{\text{alarm}} - 15 \in \{0,\dots,14\}$. A
clip that never alarms is \textbf{right-censored} at $14$: its latency is not
$14$, it is unobserved and at least $14$.

This censoring is the reason the protocol's headline metric needs a companion.
Once fewer than half the clips alarm, the median $\dt$ equals the censoring
value for \emph{every} monitor, and the pre-registered comparison returns ``no
difference'' no matter how large the true effect. We registered the median as
primary because it was specified in advance, and registered the
\textbf{restricted mean detection latency}
\begin{equation}
\RMTA \;=\; \mathbb{E}\bigl[\min(\dt, 14)\bigr] \;=\; \sum_{u=0}^{13} \widehat S(u),
\end{equation}
the area under the Kaplan--Meier curve $\widehat S$, as a secondary that stays
well defined at any detection rate. Lower is earlier.

\subsection{Tests}

\begin{itemize}
\item \textbf{Primary (as registered).} Paired Wilcoxon signed-rank on $\dt$,
probe versus baseline, one-sided (the hypothesis is directional), with censored
clips entered at $14$ --- conservative, since it understates differences. The
two-sided $p$ is always reported alongside, so the tail cannot be chosen post
hoc.
\item \textbf{Censoring-aware.} Log-rank on the Kaplan--Meier curves. Both arms
come from the same clips, violating its independence assumption, so it is
reported as corroborative only.
\item \textbf{Effect size.} Paired cluster bootstrap on the \RMTA{} difference
($\lgwNBoot$ resamples, the same resample applied to both monitors so that
pairing is preserved), giving a confidence interval in frames.
\item \textbf{Multiplicity.} Holm's step-down correction across all six
comparisons ($3$ probes $\times$ $2$ baselines), separately for the median and
\RMTA{} families.
\end{itemize}

\subsection{The decision gate}

A probe counts as a \textbf{faster detector} iff, against \emph{both} baselines, its
median $\dt$ is at least one frame lower \emph{and} the Holm-adjusted one-sided
$p$ is below $0.01$. If no probe qualifies, the protocol terminates the
expansion track and the result is written up as a negative one. The gate is
evaluated by code, not by reading tables.

\subsection{Pre-specified follow-up: what lives in the lag window}
\label{sec:mech}

If a probe does alarm before the loss, the frames in between are where the
representation has registered something the output error has not, and the
natural question is what. We pre-specify the analysis here so that it is fixed
before the lag window is known: fit linear probes to the latents $z_t$ for
$t \in [t^\star + \dt_{\text{probe}},\; t^\star + \dt_{\text{loss}})$, targeting
(i) which of the four anomalies was injected, and (ii) the disc's position and
velocity, which are the quantities the physics conditions actually change.

Two controls are part of the specification, because a linear probe on
high-dimensional latents is easy to over-read. Frames from one clip are highly
correlated, so a random train/test split lets the probe memorise clip identity
and return a near-perfect score that means nothing; all cross-validation is
therefore grouped by clip. And an accuracy is uninterpretable without a
reference, so every score is reported beside the majority-class rate and a
\emph{clip-level} label-shuffled control, which must land at chance. If the
shuffled control does not land at chance, the split is leaking and the real
score is not reported as evidence.

This analysis is contingent on the gate and is listed here as method, not
result.

\subsection{Guards against a vacuous verdict}

Two checks sit outside the hypothesis tests and exist to stop the gate from
being believed when it should not be.

\paragraph{Lift.} For every monitor we report
$\text{lift} = \text{detection rate} - \text{held-out clip false-alarm rate}$.
A monitor with lift near zero is firing at chance. If a \emph{baseline} fires at
chance, then beating it carries no information about H1 --- a probe cannot be shown to fire
\emph{sooner} than a signal that never fires --- and a passing
verdict is reported as \textsc{vacuous}. This check is not a hypothesis test and
does not alter any threshold or statistic; it only makes an existing failure
mode visible.

\paragraph{Causality, tested adversarially.} Claiming that frame $t$ cannot see
frame $t{+}1$ is cheap; testing it is not. For every registered adapter we take
a clip, invert every frame after a cut point, and assert that $z_t$,
$\mathcal{L}_0$ and all three probes are unchanged to floating-point tolerance
for $t$ at or before the cut. A separate test asserts that the $\mathcal{L}_0$
\emph{prediction} does not move when only $f_{t-1}$ and $f_t$ change, while
confirming it does move when $f_{t-2}$ changes --- so the guarantee is not
vacuously satisfied by a prediction that ignores its input. Probe streams are
additionally checked for prefix invariance: scores for frames $0..t$ must be
identical whether or not later frames exist.
